% GENERATED FILE. Edit canonical lessons, then run npm run generate:books.
% canonical-source-hash: fnv1a64:9f16eacf9e67a14f
% canonical-lessons: RU-C228-oshibka, RU-C228-pravilo, RU-C228-primer, RU-C228-znacheniye, RU-C228-imya

\chapter{Mistake, Rule, Example, Meaning, First name}
\label{ch:ru-mistake-rule-example-meaning-first-name}
\hlchaptermodality{228}

\begin{chapteropening}
\textbf{By the end of this chapter:} \emph{I can say a mistake, a rule, an example, a meaning and a first name.}

Complete the last lesson of chapter 228: I can say a mistake, a rule, an example, a meaning and a first name.
\end{chapteropening}

\section[\texorpdfstring{\ru{ошибка} (oshíbka)}{ошибка (oshíbka)}]{\ru{ошибка} (oshíbka) — a mistake}

\label{lesson:RU-C228-oshibka}



\begin{quote}
Before the new one: say the Russian for a woman, then the Russian for a boy.
\end{quote}

\subsection*{You'll want to know: \ru{ошибка}}
\textbf{\ru{ошибка}} (\emph{oshíbka}) — \textquotedblleft{}a mistake\textquotedblright{}.

\begin{cousinweb}
\textbf{\ru{ошибка}} comes from \emph{\ru{ошибиться}}, \textquotedblleft{}to make a mistake\textquotedblright{}, first \textquotedblleft{}to miss, to strike past the target\textquotedblright{}, from \emph{\ru{шибать}}, \textquotedblleft{}to strike\textquotedblright{} (Vasmer).
\end{cousinweb}

\begin{grammarlens}[title={What you've built}]
One more word for this chapter.
\end{grammarlens}

\subsection*{Your turn}
\begin{itemize}
\raggedright
  \item \emph{Say it:} \emph{oshíbka}
  \item \emph{Say it:} \emph{oshíbka}, once more
  \item \emph{Say it:} say \emph{oshíbka}
  \item \emph{Recall:} say the Russian for a woman, then the Russian for a boy, then say \emph{oshíbka} again
\end{itemize}

\begin{tcolorbox}[breakable,colback=teal!4,colframe=teal!35!black,title={Before you move on}]
What does \ru{ошибка} mean? (\textquotedblleft{}A mistake\textquotedblright{}.) Say it once more.
\end{tcolorbox}
\section[\texorpdfstring{\ru{правило} (právilo)}{правило (právilo)}]{\ru{правило} (právilo) — a rule}

\label{lesson:RU-C228-pravilo}



\begin{quote}
Before the new one: say the Russian for a boy, then the Russian for a mistake.
\end{quote}

\subsection*{You'll want to know: \ru{правило}}
\textbf{\ru{правило}} (\emph{právilo}) — \textquotedblleft{}a rule\textquotedblright{}.

\begin{cousinweb}
\textbf{\ru{правило}} is \emph{\ru{править}}, \textquotedblleft{}to set straight, to rule\textquotedblright{}, with \emph{-\ru{ло}}, a suffix for a tool: \textquotedblleft{}what sets things straight\textquotedblright{}. \emph{\ru{правый}} is \textquotedblleft{}straight, right\textquotedblright{}.
\end{cousinweb}

\begin{grammarlens}[title={What you've built}]
One more word for this chapter.
\end{grammarlens}

\subsection*{Your turn}
\begin{itemize}
\raggedright
  \item \emph{Say it:} \emph{právilo}
  \item \emph{Say it:} \emph{právilo}, once more
  \item \emph{Say it:} say \emph{právilo}
  \item \emph{Recall:} say the Russian for a boy, then the Russian for a mistake, then say \emph{právilo} again
\end{itemize}

\begin{tcolorbox}[breakable,colback=teal!4,colframe=teal!35!black,title={Before you move on}]
What does \ru{правило} mean? (\textquotedblleft{}A rule\textquotedblright{}.) Say it once more.
\end{tcolorbox}
\section[\texorpdfstring{\ru{пример} (primér)}{пример (primér)}]{\ru{пример} (primér) — an example}

\label{lesson:RU-C228-primer}



\begin{quote}
Before the new one: say the Russian for a mistake, then the Russian for a rule.
\end{quote}

\subsection*{You'll want to know: \ru{пример}}
\textbf{\ru{пример}} (\emph{primér}) — \textquotedblleft{}an example\textquotedblright{}.

\begin{cousinweb}
\textbf{\ru{пример}} is \emph{\ru{при}-}, \textquotedblleft{}by, against\textquotedblright{}, with \emph{\ru{мер}}, from \emph{\ru{мера}}, \textquotedblleft{}a measure\textquotedblright{}: something to measure against. Vasmer relates \emph{\ru{мера}} to Latin \emph{mētior}, \textquotedblleft{}I measure\textquotedblright{}.
\end{cousinweb}

\begin{grammarlens}[title={What you've built}]
One more word for this chapter.
\end{grammarlens}

\subsection*{Your turn}
\begin{itemize}
\raggedright
  \item \emph{Say it:} \emph{primér}
  \item \emph{Say it:} \emph{primér}, once more
  \item \emph{Say it:} say \emph{primér}
  \item \emph{Recall:} say the Russian for a mistake, then the Russian for a rule, then say \emph{primér} again
\end{itemize}

\begin{tcolorbox}[breakable,colback=teal!4,colframe=teal!35!black,title={Before you move on}]
What does \ru{пример} mean? (\textquotedblleft{}An example\textquotedblright{}.) Say it once more.
\end{tcolorbox}
\section[\texorpdfstring{\ru{значение} (znachéniye)}{значение (znachéniye)}]{\ru{значение} (znachéniye) — a meaning}

\label{lesson:RU-C228-znacheniye}



\begin{quote}
Before the new one: say the Russian for a rule, then the Russian for an example.
\end{quote}

\subsection*{You'll want to know: \ru{значение}}
\textbf{\ru{значение}} (\emph{znachéniye}) — \textquotedblleft{}a meaning\textquotedblright{}.

\begin{cousinweb}
\textbf{\ru{значение}} comes from \emph{\ru{значить}}, \textquotedblleft{}to mean\textquotedblright{}, from \emph{\ru{знак}}, \textquotedblleft{}a sign\textquotedblright{}, which goes with \emph{\ru{знать}}, \textquotedblleft{}to know\textquotedblright{} (Vasmer).
\end{cousinweb}

\begin{grammarlens}[title={What you've built}]
One more word for this chapter.
\end{grammarlens}

\subsection*{Your turn}
\begin{itemize}
\raggedright
  \item \emph{Say it:} \emph{znachéniye}
  \item \emph{Say it:} \emph{znachéniye}, once more
  \item \emph{Say it:} say \emph{znachéniye}
  \item \emph{Recall:} say the Russian for a rule, then the Russian for an example, then say \emph{znachéniye} again
\end{itemize}

\begin{tcolorbox}[breakable,colback=teal!4,colframe=teal!35!black,title={Before you move on}]
What does \ru{значение} mean? (\textquotedblleft{}A meaning\textquotedblright{}.) Say it once more.
\end{tcolorbox}
\section[\texorpdfstring{\ru{имя} (ímya)}{имя (ímya)}]{\ru{имя} (ímya) — a first name}

\label{lesson:RU-C228-imya}



\begin{quote}
Before the new one: say the Russian for an example, then the Russian for a meaning.
\end{quote}

\subsection*{You'll want to know: \ru{имя}}
\textbf{\ru{имя}} (\emph{ímya}) — \textquotedblleft{}a first name\textquotedblright{}.

\begin{cousinweb}
\textbf{\ru{имя}} $\leftarrow$ Proto-Slavic \emph{*jĭmę} — inherited. Vasmer lists Old Prussian \emph{emmens}, Latin \emph{nōmen}, Gothic \emph{namō} and Old High German \emph{-nuomen}: English \emph{name} is a cousin, from PIE \emph{*h\textsubscript{1}nómṇ}.
\end{cousinweb}

\begin{grammarlens}[title={What you've built}]
One more word for this chapter.
\end{grammarlens}

\subsection*{Your turn}
\begin{itemize}
\raggedright
  \item \emph{Say it:} \emph{ímya}
  \item \emph{Say it:} \emph{ímya}, once more
  \item \emph{Say it:} say \emph{ímya}
  \item \emph{Recall:} say the Russian for an example, then the Russian for a meaning, then say \emph{ímya} again
\end{itemize}

\begin{tcolorbox}[breakable,colback=teal!4,colframe=teal!35!black,title={Before you move on}]
What does \ru{имя} mean? (\textquotedblleft{}A first name\textquotedblright{}.) Say it once more.
\end{tcolorbox}
